% TOP_PROBLEM: {"id": 30007461, "title": "Integer multiplication below n log n", "category_id": 15, "category": {"id": 15, "name": "computer_science", "display_name": "Computer Science", "description": "Computational complexity, algorithms, and theoretical CS.", "slug": "computer-science", "order_index": 15, "created_at": "2026-07-31T15:26:25.671Z"}, "status": "open"}
% FIRST_POSTED: 2026-10-06
% ATTEMPT_STATUS: solved
% OPENAI_CLAIM: {"schema_version":1,"source_repo":"https://github.com/openai/math","source_commit":"adc7f1241b42e322a6451854ab7e4b4c146bf78a","release_date":"2026-10-06","match":"full","resolution":"proved","independently_reviewed":false,"second_pass":"agreed","families":[{"family":"109","title":"Integer multiplication below $`n\\log n`$","subject":"Theoretical computer science","manuscripts":[{"title":"Integer multiplication below nlog n","pdf_path":"preprints/Integer-multiplication-below-n-log-n-September-23-2026/paper.pdf","readme_path":"preprints/Integer-multiplication-below-n-log-n-September-23-2026/README.md","date":"2026-09-23","theorem_ref":"thm:main (sections/00-introduction.tex:21)","statement_sources":[{"path":"preprints/Integer-multiplication-below-n-log-n-September-23-2026/build/sections/00-introduction.tex","line":21,"label":"thm:main"}]}],"lean":null,"reasoning_trace":null}]}
\documentclass[11pt]{article}
\usepackage[margin=25mm]{geometry}
\usepackage{fontspec}
\setmainfont{Latin Modern Roman}
\usepackage{amsmath,amssymb,mathtools}
\usepackage{unicode-math}
\setmathfont{Latin Modern Math}
\usepackage{enumitem}
\usepackage{xurl,hyperref,bookmark}
\hypersetup{colorlinks=true,urlcolor=blue}
\setcounter{secnumdepth}{0}
\setlength{\parindent}{0pt}
\setlength{\parskip}{5pt}
\setlength{\emergencystretch}{3em}
\begin{document}
\section*{Integer multiplication below n log n}

\subsection{OpenAI's claim}
OpenAI presents the following claims relevant to this question, with manuscript numbers following paper order in Sources. In manuscript 1 (family 109; 2026-09-23), the asserted result is that exact multiplication on a fixed multitape finite-alphabet Turing machine takes $O(n(\log n)^{1-2^{-182}})$ worst-case time.

\subsection{Mathematical statement of the claimed result}
\paragraph{Manuscript 1.}
For every $n\ge1$ and two binary words of length $n$, the claim supplies one fixed deterministic Turing machine, with fixed finite alphabet and fixed number of one-dimensional tapes, that outputs the exact product of their nonnegative integer values in worst-case time
\[
 O(n(\lg n)^{1-\kappa}),\qquad\kappa=2^{-182}.
\]
The convention is $\lg n=\max\{\lceil\log_2n\rceil,1\}$. In the same time scale, fixed machines produce quotient and remainder for a nonzero divisor and the integer square root, with padded binary outputs of length $n$ allowed. These are bit-operation bounds, not unit-cost real arithmetic.

\subsection{Outline of the argument}
For manuscript 1 (family 109; 2026-09-23), two finite linear networks exchange data banks while restoring all auxiliary values. Subspace labels make representation changes cheap on array wires, and the rank saving survives recursive transforms; exact reconstruction and tape-stream accounting give the bit-time improvement.

\subsection{Scope relative to this problem}
The result supplies the conclusion in this problem with its displayed hypotheses and quantifiers. The broader statements and applications in the manuscript are kept separate from the particular question above. The supporting result is manuscript 1 (family 109; 2026-09-23).

\subsection{Formal verification}
No Lean formalization is listed at this snapshot for manuscript 1 (family 109; 2026-09-23).

This site has not executed the Lean code. Described statement coverage is not an independent proof check.

\subsection{Status caveat}
The manuscripts were produced by an internal OpenAI model. OpenAI cautions that results without formalization can contain mistakes. This site has not independently checked the proofs; the label records the scope of OpenAI's claim.

\subsection{Sources}
\begin{itemize}
\item OpenAI, \emph{Integer multiplication below nlog n}, 2026-09-23; citation key \texttt{OAI:Integer-multiplication-below-n-log-n-September-23-2026}. \href{https://github.com/openai/math/blob/adc7f1241b42e322a6451854ab7e4b4c146bf78a/preprints/Integer-multiplication-below-n-log-n-September-23-2026/paper.pdf}{PDF (GitHub)}; \href{https://github.com/openai/math/raw/adc7f1241b42e322a6451854ab7e4b4c146bf78a/preprints/Integer-multiplication-below-n-log-n-September-23-2026/paper.pdf}{PDF (direct)}; \href{https://github.com/openai/math/blob/adc7f1241b42e322a6451854ab7e4b4c146bf78a/preprints/Integer-multiplication-below-n-log-n-September-23-2026/README.md}{Manuscript page and citation}; \href{https://github.com/openai/math/blob/main/preprints/Integer-multiplication-below-n-log-n-September-23-2026/paper.pdf}{Latest version}; Statement source: \href{https://github.com/openai/math/blob/adc7f1241b42e322a6451854ab7e4b4c146bf78a/preprints/Integer-multiplication-below-n-log-n-September-23-2026/build/sections/00-introduction.tex\#L21}{thm:main}.
\item \href{https://github.com/openai/math/blob/adc7f1241b42e322a6451854ab7e4b4c146bf78a/CONTENTS.md}{OpenAI catalogue entry}.
\end{itemize}
\end{document}
